% AA-31 | Iteracja XVI; szczegolowy audyt historyczny IV
\chapter{Wielomiany symetryczne i wzory Viète'a}\label{aa:ch:symmetric-polynomials}
\section{Działanie grupy symetrycznej}

W obliczeniach z pierwiastkami wielomianu często nie znamy samych
pierwiastków, lecz znamy ich sumę i iloczyn. Zamiana dwóch
pierwiastków nie zmienia tych danych. Powstaje pytanie, czy każdą
wielomianową zależność niezmienniczą przy zamianie pierwiastków
można zapisać wyłącznie za pomocą pewnych podstawowych sum
symetrycznych. Odpowiedź jest twierdząca nawet wtedy, gdy
współczynniki należą do pierścienia z dzielnikami zera.

\begin{definicja}\label{aa:def:symmetric-poly}
Niech $A$ będzie przemiennym pierścieniem z jedynką
i niech $n\ge1$ będzie liczbą całkowitą. Grupa $S_n$ składa się
ze wszystkich bijekcji zbioru $\{1,\ldots,n\}$ na siebie.
Wielomian $f\in A[X_1,\ldots,X_n]$ jest
\emph{symetryczny}, jeśli dla każdej permutacji $\sigma\in S_n$
zachodzi
\[
 f(X_{\sigma(1)},\ldots,X_{\sigma(n)})=f(X_1,\ldots,X_n).
\]
Zbiór wszystkich takich wielomianów, z działaniami odziedziczonymi
z $A[X_1,\ldots,X_n]$, oznaczamy przez
$A[X_1,\ldots,X_n]^{S_n}$.
Dla $0\le k\le n$ oznaczamy przez $e_k$ $k$-ty
\emph{elementarny wielomian symetryczny}:
\[
 e_0=1,\qquad
 e_k=\sum_{1\le i_1<\cdots<i_k\le n}X_{i_1}\cdots X_{i_k}
 \quad(1\le k\le n).
\]
Każdy składnik sumy odpowiada wyborowi $k$ \emph{różnych}
indeksów; suma nie zawiera powtórzeń tego samego podzbioru.
Każda permutacja zmiennych permutuje rodzinę takich podzbiorów,
więc $e_k\in A[X_1,\ldots,X_n]^{S_n}$ dla $0\le k\le n$.
\end{definicja}

\begin{przyklad}[Trzy zmienne]
Dla $n=3$ otrzymujemy
\[
 e_1=X_1+X_2+X_3,\qquad
 e_2=X_1X_2+X_1X_3+X_2X_3,\qquad
 e_3=X_1X_2X_3.
\]
Po zamianie $X_1$ z $X_3$ składniki pierwszej sumy jedynie
zmieniają kolejność, podobnie jak trzy składniki drugiej sumy.
Iloczyn $e_3$ również pozostaje bez zmian. Tak samo działa
każda permutacja, więc $e_k$ należy do pierścienia wielomianów
symetrycznych.
\end{przyklad}

\begin{twierdzenie}[Podstawowe twierdzenie o wielomianach symetrycznych]\label{aa:tw:fundamental-symmetric}
Dla każdego przemiennego pierścienia $A$ z jedynką
i każdego $n\ge1$ homomorfizm pierścieni
\[
 \Phi:A[T_1,\ldots,T_n]\longrightarrow
 A[X_1,\ldots,X_n]^{S_n},\qquad
 \Phi(T_k)=e_k\quad(1\le k\le n),\quad\Phi(a)=a\quad(a\in A),
\]
jest izomorfizmem. W szczególności każdy wielomian symetryczny
$f$ ma \emph{jednoznaczne} przedstawienie
$f=P(e_1,\ldots,e_n)$ dla pewnego
$P\in A[T_1,\ldots,T_n]$.
\end{twierdzenie}

\begin{proof}
Jeżeli $1_A=0$, to $A$ jest pierścieniem zerowym.
Wówczas oba pierścienie występujące w twierdzeniu
są jednopunktowe i $\Phi$ jest izomorfizmem.
W dalszej części przyjmujemy $1_A\ne0$.

\textbf{Krok 1: porządek i pierwszy jednomian.}
Jednomiany $X^a=X_1^{a_1}\cdots X_n^{a_n}$, gdzie
$a=(a_1,\ldots,a_n)\in\mathbb N_0^n$, porządkujemy
leksykograficznie przy $X_1>\cdots>X_n$.
Oznacza to, że $X^a>X^b$, jeśli w pierwszej współrzędnej,
w której $a$ i $b$ się różnią, mamy $a_j>b_j$.
Każdy niezerowy wielomian ma skończenie wiele wyrazów
o niezerowych współczynnikach, a zatem ma największy wyraz,
zwany wiodącym. Jeśli $f$ jest symetryczny i jego wyraz
wiodący jest równy $cX_1^{a_1}\cdots X_n^{a_n}$, gdzie
$c\in A\setminus\{0\}$, to $a_1\ge\cdots\ge a_n$.
Rzeczywiście, gdyby dla pewnego $i<j$ zachodziło $a_i<a_j$,
permutacja zamieniająca $X_i$ z $X_j$ wymuszałaby wystąpienie
w $f$ wyrazu o tej samej niezerowej wartości współczynnika
$c$ i o większym jednomianie. To przeczy wyborowi wyrazu
wiodącego. Wniosek nie wymaga dzielenia przez żaden element $A$.

\textbf{Krok 2: eliminacja wyrazu wiodącego.}
Wyrazem wiodącym $e_k$ jest $X_1\cdots X_k$ z
współczynnikiem $1_A$. Dla $a_1\ge\cdots\ge a_n$ kładziemy
$a_{n+1}=0$ i rozpatrujemy
\[
 M_a=e_1^{a_1-a_2}e_2^{a_2-a_3}\cdots e_n^{a_n-a_{n+1}}.
\]
Wykładniki są nieujemne. Mnożenie wyrazów wiodących
daje dokładnie $X_1^{a_1}\cdots X_n^{a_n}$, ponieważ
wykładnik $X_j$ wynosi
$(a_j-a_{j+1})+\cdots+(a_n-a_{n+1})=a_j$.
Ponieważ współczynniki wiodące wszystkich $e_k$ są
równe $1_A$, współczynnik tego wyrazu w $M_a$ również
jest $1_A$, niezależnie od istnienia dzielników zera w $A$.
Wielomian $f-cM_a$ jest symetryczny, a jego wyraz wiodący
jest ściśle mniejszy od dotychczasowego (chyba że otrzymaliśmy
wielomian zerowy).

\textbf{Krok 3: skończoność algorytmu.}
Wszystkie jednomiany $M_a$ są jednorodne stopnia
$a_1+\cdots+a_n$. W szczególności żaden nowy wyraz
powstały przy odejmowaniu $cM_a$ nie ma stopnia
większego niż początkowy $d=\deg f$.
W $n$ zmiennych istnieje tylko skończenie wiele jednomianów
stopnia całkowitego nie większego niż $d$: są one
indeksowane przez wektory $b\in\mathbb N_0^n$ o
$b_1+\cdots+b_n\le d$. Przy każdym niezakończonym
kroku jednomian wiodący ściśle maleje w tym skończonym
zbiorze. Algorytm zatem się kończy i wyraża $f$ jako
wielomian w $e_1,\ldots,e_n$. Jest to surjektywność $\Phi$.

\textbf{Krok 4: jednoznaczność przedstawienia.}
Niech $P\in A[T_1,\ldots,T_n]$ będzie niezerowy.
Dla każdego jednomianu $T^b$ jego obraz ma współczynnik
wiodący $1_A$, ponieważ wielomiany $e_k$ są moniczne
w ustalonym porządku leksykograficznym.
Jednomianowi $T^b=T_1^{b_1}\cdots T_n^{b_n}$ odpowiada
jednomian wiodący
\[
 X_1^{b_1+\cdots+b_n}X_2^{b_2+\cdots+b_n}
 \cdots X_n^{b_n}
\]
po podstawieniu $T_k=e_k$. Mapa wykładników
\[
 (b_1,\ldots,b_n)\longmapsto
 (b_1+\cdots+b_n,b_2+\cdots+b_n,\ldots,b_n)
\]
jest różnowartościowa: z prawej strony odzyskujemy
$b_j$ jako różnicę $j$-tej i $(j+1)$-szej współrzędnej
(przyjmując ostatnią dodatkową współrzędną równą $0$).
Spośród niezerowych wyrazów $P$ wybierzmy ten, którego
obraz ma największy jednomian wiodący. Inne wyrazy
mają mniejsze jednomiany wiodące, a każdy niewiodący
składnik ich rozwinięcia jest tym bardziej mniejszy.
Zatem wybrany największy jednomian nie może się skasować:
pochodzi tylko z wybranego wyrazu $P$, a jego współczynnik
pozostaje ten sam, niezerowy w $A$.
Otrzymujemy $\Phi(P)\ne0$, co dowodzi iniektywności.
\end{proof}

\begin{przyklad}[Algorytm dla dwóch zmiennych]\label{aa:prz:symmetric-algorithm}
Weźmy $f=X_1^2+X_2^2+X_1X_2$ nad $\mathbb Z$.
Wyrazem wiodącym jest $X_1^2$, więc odejmujemy $e_1^2$:
\[
 f-e_1^2
 =X_1^2+X_2^2+X_1X_2-(X_1^2+2X_1X_2+X_2^2)
 =-X_1X_2=-e_2.
\]
Oznacza to $f=e_1^2-e_2$. Nie zgadywaliśmy współczynników:
wyznaczył je algorytm eliminowania wyrazu wiodącego.
\end{przyklad}

\begin{przyklad}[Pierścień z dzielnikami zera]
Dla $A=\mathbb Z/4\mathbb Z$ i dwóch zmiennych
wielomian $X_1^2+X_2^2$ nadal spełnia
$X_1^2+X_2^2=e_1^2-2e_2$.
Nie można tutaj skrócić dowodu przez przejście do
ciała ułamków $A$: pierścień $A$ ma niezerowe dzielniki zera,
więc nie posiada ciała ułamków w zwykłym znaczeniu.
Dowód powyższego twierdzenia działa bez takiego przejścia,
ponieważ korzysta jedynie ze współczynników wiodących $1_A$.
\end{przyklad}

\begin{przyklad}[Eliminacja wyrazu wiodącego w trzech zmiennych]\label{aa:prz:symmetric-three-elimination}
Niech $A=\mathbb Z$ i
$f=X_1^2X_2+X_1^2X_3+X_2^2X_1+X_2^2X_3+X_3^2X_1+X_3^2X_2$.
Jest to suma wszystkich jednomianów $X_i^2X_j$ z $i\ne j$;
permutacja zmiennych jedynie zmienia kolejność sześciu składników.
Wyraz wiodący w porządku leksykograficznym jest $X_1^2X_2$.
Odpowiada mu wektor wykładników $(2,1,0)$, a więc
$M_a=e_1^{2-1}e_2^{1-0}=e_1e_2$.
Rozwijając iloczyn, otrzymujemy
\[
 e_1e_2=f+3X_1X_2X_3=f+3e_3.
\]
Współczynnik $3$ pochodzi stąd, że jednomian $X_1X_2X_3$
można uzyskać na trzy sposoby: wybierając w $e_1$ kolejno
$X_1$, $X_2$ lub $X_3$, a w $e_2$ iloczyn pozostałych dwóch
zmiennych. Stąd $f=e_1e_2-3e_3$.
Tożsamość pozostaje poprawna po przejściu do dowolnego
przemiennego pierścienia z jedynką, także gdy jego charakterystyka
dzieli $3$.
\end{przyklad}

\section{Pierwiastki i współczynniki}

Każdy wielomian moniczny dodatniego stopnia nad ciałem
rozkłada się w swoim ciele rozkładu na czynniki liniowe.
Zestawienie jego współczynników
z wielomianami symetrycznymi wyliczanymi na pierwiastkach wyjaśnia,
dlaczego nie musimy znać poszczególnych pierwiastków,
aby obliczać niektóre wyrażenia z ich udziałem.

\begin{twierdzenie}[Wzory Viète'a]\label{aa:tw:vieta-general}
Niech $K$ będzie ciałem, $n\ge1$ i
$f(X)=X^n+c_1X^{n-1}+\cdots+c_n\in K[X]$.
Niech $L/K$ będzie ciałem rozkładu $f$, a
$\alpha_1,\ldots,\alpha_n\in L$ jego pierwiastkami,
wypisanymi z krotnościami algebraicznymi.
Wówczas dla każdej liczby $k\in\{1,\ldots,n\}$ zachodzi
$e_k(\alpha_1,\ldots,\alpha_n)=(-1)^kc_k$.
\end{twierdzenie}
\begin{proof}
W $L[X]$ zachodzi
$f(X)=\prod_{j=1}^n(X-\alpha_j)$, gdzie każdy pierwiastek
jest uwzględniony tyle razy, ile wynosi jego krotność.
W rozwinięciu tego iloczynu wyraz stopnia $n-k$ w $X$
otrzymujemy, wybierając $k$ spośród $n$ czynników,
z których bierzemy $-\alpha_j$, a z pozostałych
bierzemy $X$. Zbiór wyborów jest dokładnie zbiorem
$k$-elementowych podzbiorów $\{1,\ldots,n\}$.
Współczynnik przy $X^{n-k}$ jest zatem równy
$(-1)^ke_k(\alpha_1,\ldots,\alpha_n)$.
Porównanie go ze współczynnikiem $c_k$ po lewej stronie
daje $c_k=(-1)^ke_k$. Mnożąc przez $(-1)^k$
(dla dowolnej charakterystyki) otrzymujemy tezę.
\end{proof}

\begin{przyklad}[Suma kwadratów pierwiastków]\label{aa:prz:symmetric-squares}
Dla monicznego wielomianu
$f(X)=X^3-2X^2+3X-4\in\mathbb Q[X]$ mamy ze wzorów
Viète'a $e_1(\alpha_1,\alpha_2,\alpha_3)=2$ oraz
$e_2(\alpha_1,\alpha_2,\alpha_3)=3$.
Rozwijając kwadrat sumy, otrzymujemy
\[
 (\alpha_1+\alpha_2+\alpha_3)^2
 =\alpha_1^2+\alpha_2^2+\alpha_3^2
  +2(\alpha_1\alpha_2+\alpha_1\alpha_3+\alpha_2\alpha_3).
\]
Stąd suma kwadratów wynosi $e_1^2-2e_2=4-6=-2$.
Obliczenie nie wymaga wyznaczania żadnego z trzech pierwiastków.
\end{przyklad}

\begin{przyklad}[Suma sześcianów]
Niech $f(X)=X^3+c_1X^2+c_2X+c_3$ będzie moniczny
nad dowolnym ciałem $K$. Tożsamość w pierścieniu
$\mathbb Z[U,V,W]$ daje
\[
 U^3+V^3+W^3=(U+V+W)^3
 -3(U+V+W)(UV+UW+VW)+3UVW.
\]
Dla pierwiastków $f$, po podstawieniu
$e_1=-c_1$, $e_2=c_2$, $e_3=-c_3$, otrzymujemy
$\sum_j\alpha_j^3=-c_1^3+3c_1c_2-3c_3$.
Ponieważ wykorzystana tożsamość ma współczynniki
całkowite, ten sam wzór jest poprawny również
w charakterystykach $2$ i $3$ (odczytując odpowiednie
współczynniki w $K$).
\end{przyklad}

\begin{przyklad}[Kontrprzykład: symetria nie oznacza liniowości w $e_k$]\label{aa:prz:symmetric-nonlinear}
Wielomian $X_1^2+X_2^2$ jest symetryczny, lecz nie ma
postaci $a e_1+b e_2$ dla $a,b\in\mathbb Z$:
$e_1=X_1+X_2$ ma tylko wyrazy stopnia $1$, a
$e_2=X_1X_2$ jest jednomianem stopnia $2$ bez
wyrazu $X_1^2$. Porównanie współczynnika przy
$X_1^2$ prowadzi do sprzeczności $1=0$.
Poprawne przedstawienie jest nieliniowe:
$X_1^2+X_2^2=e_1^2-2e_2$.
\end{przyklad}

\begin{uwaga}[Ostrzeżenie w dodatniej charakterystyce]
W charakterystyce $2$ tożsamość
$X_1^2+X_2^2=e_1^2-2e_2$ upraszcza się do
$X_1^2+X_2^2=e_1^2$.
Nie wolno zatem próbować odzyskać $e_2$ z sumy
kwadratów przez dzielenie przez $2$: liczba $2$
jest zerem w ciele charakterystyki $2$.
Nie przeczy to podstawowemu twierdzeniu o wielomianach
symetrycznych: jego generatorami są nadal $e_1,e_2$,
a nie jedynie suma kwadratów.
\end{uwaga}

\begin{uwaga}\label{aa:uw:symmetric-galois}
Niech $f\in K[X]$ będzie rozdzielczy i niech $L/K$
będzie jego ciałem rozkładu. Grupa
$G=\operatorname{Gal}(L/K)$ permutuje pierwiastki $f$.
Każdy wielomian symetryczny o współczynnikach w $K$,
obliczony na tych pierwiastkach, należy do $K$:
wynika to bezpośrednio z Twierdzenia~\ref{aa:tw:fundamental-symmetric}
i wzorów Viète'a (Twierdzenie~\ref{aa:tw:vieta-general}),
bez korzystania z odpowiedniości Galois.
Nie trzeba też zakładać rozdzielczości dla tej ostatniej
konkluzji, jeżeli pierwiastki liczymy z krotnościami.
\end{uwaga}
